\chapter{Technische Grundlagen}
\abschnittsautor{S. Kaufmann, M. Raichle, P. Riedl}

Hier steht alles, was zum Verständnis der praktischen Umsetzung nötig ist --
und nur das. Jede Technologie, die ihr später verwendet, wird hier einmal
erklärt und belegt.

\section{Teilbereich A}

Beim ersten Auftreten schreibt ihr eine Abkürzung aus: \ac{API}. Danach
genügt die Kurzform: \ac{API}. Alle Abkürzungen sammelt ihr in
\texttt{diplomarbeit.tex} im Abkürzungsverzeichnis.

\subsection{Ein Unterabschnitt}

Zitiert wird im IEEE-Stil mit einer Nummer in eckigen
Klammern~\cite{beispiel-artikel}. Mehrere Quellen auf einmal gehen
auch~\cite{beispiel-buch,beispiel-online}.

\subsubsection{Eine Ebene tiefer}

Die vierte Ebene könnt ihr verwenden, aber sparsam. Sie wird nummeriert,
steht aber nicht mehr im Inhaltsverzeichnis. Tiefer gliedert ihr nicht.

\section{Teilbereich B}

Erklärt Technologien in eurer eigenen Formulierung. Wörtlich übernommene
Passagen gehören in Anführungszeichen und mit Seitenangabe belegt, sonst gilt
das als Plagiat.



\newpage
\section{Teilbereich C: Sprachmodellgestützte Sicherheitsanalyse}
\label{subsec:teilbereich-c}

\subsection{Zielsetzung und technische Anforderungen}

Ziel dieses Teilbereichs ist die Planung einer
sprachmodellgestützten Sicherheitsvorfallsbewertung.
Das System soll technische Ereignisdaten aufbereiten,
deren sicherheitsrelevante Bedeutung erläutern und geeignete
technische Maßnahmen vorschlagen. Zusätzlich
sollen strukturierte Regelvorschläge erzeugt werden, die im weiteren
Verlauf geprüft und kontrolliert vom Nutzer freigegeben werden können.

Ein wesentlicher Anwendungsfall ist die Unterstützung bei dera
Untersuchung von Ereignissen, die beispielsweise durch ungewöhnliche
Netzwerkverbindungen, fehlgeschlagene Anmeldeversuche oder andere
sicherheitsrelevante Aktivitäten ausgelöst werden. Das Sprachmodell
soll die verfügbaren Informationen zusammenfassen, mögliche
Ursachen erläutern und die vorgeschlagenen Maßnahmen nachvollziehbar
begründen.

Dabei muss zwischen tatsächlich beobachteten Informationen und
Interpretationen des Modells unterschieden werden. Ein einzelner
fehlgeschlagener Anmeldeversuch belegt beispielsweise noch keinen
erfolgreichen Angriff. Reichen die vorhandenen Daten für eine
verlässliche Bewertung nicht aus, soll das Modell auf fehlende
Informationen hinweisen, anstatt unbelegte Schlussfolgerungen zu
erzeugen.

Neben der fachlichen Qualität sind die Einhaltung vorgegebener
Ausgabeformate, die lokale Ausführbarkeit und die nachvollziehbare
Verarbeitung der Ergebnisse wichtige Anforderungen. Die erzeugten
Daten müssen sich in der Anwendung weiterverarbeiten und den
zugehörigen Sicherheitsvorfällen eindeutig zuordnen lassen.

\subsection{Auswahl des Sprachmodells}

Als Ausgangspunkt wird das bereits untersuchte lokale Qwen-Modell
mit sieben Milliarden Parametern verwendet. Dieses Modell dient als
Referenz für die weitere Bewertung, da seine Ergebnisse mit
alternativen Modellen unter vergleichbaren Bedingungen untersucht
werden können.

Als Vergleichskandidat wird Qwen3-8B vorgesehen. Die Modellfamilie
wurde unter anderem hinsichtlich Schlussfolgerungsfähigkeit,
Anweisungsbefolgung und Programmieraufgaben weiterentwickelt.
Diese Eigenschaften sind für die Analyse komplexerer
Sicherheitsvorfälle und die Erzeugung strukturierter Vorschläge
von Interesse.

Zusätzlich kann Llama 3.1 8B Instruct als Vertreter einer anderen
Modellfamilie berücksichtigt werden. Dadurch lässt sich untersuchen,
ob Unterschiede zwischen den Modellen bei den konkreten Aufgaben
des Projekts tatsächlich messbar sind.

Die Auswahl erfolgt nicht ausschließlich anhand der Modellgröße.
Entscheidend sind die fachliche Korrektheit der Analysen, die
Nachvollziehbarkeit der Begründungen, der Umgang mit unvollständigen
Informationen und die Einhaltung des vorgegebenen Ausgabeformats.
Ebenso werden die benötigte Rechenleistung und die Antwortzeit
berücksichtigt.

Die offiziellen Modellbeschreibungen dienen zunächst als Grundlage
für die Vorauswahl. Eine endgültige Entscheidung soll erst nach
vergleichbaren Tests mit einheitlichen Eingabedaten getroffen
werden.\footnote{Vgl. Qwen, \url{https://huggingface.co/Qwen/Qwen3-8B}.}
\footnote{Vgl. Meta, \url{https://huggingface.co/meta-llama/Llama-3.1-8B-Instruct}.}

\subsection{Technologische Grundlagen des Frontends}

Für die Entwicklung der Benutzeroberfläche wird eine Kombination
aus React, TypeScript und Vite vorgesehen. React ermöglicht die
Aufteilung des Dashboards in wiederverwendbare Komponenten.
Dadurch können beispielsweise Ereignislisten, Detailansichten,
Risikokennzahlen und Dialoge zur Freigabe von Regelvorschlägen
getrennt entwickelt und später gemeinsam eingesetzt werden.

TypeScript ergänzt JavaScript um ein statisches Typsystem. Damit
können die Datenstrukturen der Anwendung, etwa Sicherheitsereignisse,
Modellanalysen und Regelvorschläge, ausdrücklich definiert werden.
Fehler bei der Verwendung dieser Daten lassen sich dadurch bereits
während der Entwicklung teilweise erkennen. Die Validierung der
tatsächlich übertragenen Daten bleibt jedoch Aufgabe des Backends.

Vite wird als Entwicklungs- und Build-System vorgesehen. Es soll
eine einfache Entwicklungsumgebung für das React-Frontend
bereitstellen und die Erstellung der späteren Anwendung ermöglichen.
Da die Benutzeroberfläche hauptsächlich mit einer separaten
Backend-API kommunizieren soll, ist zunächst kein zusätzliches
Full-Stack-Framework erforderlich.

Für die Gestaltung werden Tailwind CSS und shadcn/ui vorgesehen.
Tailwind CSS unterstützt eine konsistente Gestaltung durch
wiederverwendbare Utility-Klassen. shadcn/ui stellt anpassbare
Komponenten für typische Interaktionen bereit, darunter Formulare,
Dialoge, Registerkarten und Schaltflächen. Die Komponenten können
an die Anforderungen des Sicherheitsdashboards angepasst werden.

Die Gestaltung soll eine übersichtliche Darstellung der Ereignisse
und ihrer Bearbeitungszustände ermöglichen. Kritische Informationen
müssen schnell erkennbar sein, ohne dass die Benutzeroberfläche
durch zu viele Farben oder dekorative Elemente überladen wird.

\footnote{Vgl. React, \url{https://react.dev/learn/typescript};
Vite, \url{https://vite.dev/guide/};
Tailwind CSS, \url{https://tailwindcss.com/docs/styling-with-utility-classes};
shadcn/ui, \url{https://ui.shadcn.com/docs}.}

\subsection{Backend und Kommunikation mit dem Sprachmodell}

Als Backend-Technologie wird Python mit FastAPI vorgesehen.
Das Backend übernimmt die Verarbeitung der Anfragen, die
Anbindung des lokalen Sprachmodells und die Validierung der
zurückgegebenen Daten. Die Trennung zwischen Frontend und
Backend ermöglicht es, die Benutzeroberfläche und die
Analysekomponente unabhängig voneinander weiterzuentwickeln.

Die Kommunikation erfolgt über definierte API-Endpunkte.
Das Frontend kann beispielsweise Sicherheitsvorfälle abrufen,
eine Analyse anfordern oder den Prüfstatus eines Regelvorschlags
anzeigen. Die fachliche Verarbeitung und die Berechtigungsprüfung
erfolgen dabei im Backend.

Zur Beschreibung und Validierung der übertragenen Daten werden
Pydantic-Datenmodelle vorgesehen. Diese legen fest, welche Felder
vorhanden sein müssen und welche Datentypen beziehungsweise
Wertebereiche zulässig sind. Dadurch können fehlerhafte oder
unvollständige Anfragen und Antworten gezielt erkannt werden.

Für die lokale Ausführung des Sprachmodells wird Ollama als
mögliche Laufzeitumgebung berücksichtigt. Die konkrete
Modellversion und Konfiguration werden im Zuge der Implementierung
festgelegt. Die Modellanbindung soll über eine klar definierte
Schnittstelle erfolgen, damit das Modell bei Bedarf ausgetauscht
werden kann, ohne die gesamte Anwendung neu entwickeln zu müssen.

\footnote{Vgl. FastAPI, \url{https://fastapi.tiangolo.com/};
Ollama, \url{https://docs.ollama.com/}.}

\subsection{Strukturierte Ausgaben und Datenhaltung}

Für die Weiterverarbeitung der Modellantworten wird ein
verbindliches JSON-Schema vorgesehen. Dieses legt fest, welche
Informationen eine Analyse enthalten muss. Dazu gehören
beispielsweise eine Zusammenfassung des Ereignisses, die
Begründung der Risikobewertung, empfohlene Maßnahmen und
strukturierte Angaben zu einem möglichen Regelvorschlag.

Die Verwendung eines festen Schemas erleichtert die Verarbeitung
der Modellausgaben durch das Backend und die Darstellung im
Frontend. Ein syntaktisch gültiges JSON-Dokument ist allerdings
noch kein Nachweis für die fachliche Richtigkeit seiner Inhalte.
Deshalb müssen die Ergebnisse zusätzlich anhand definierter
Validierungsregeln geprüft werden.

Besonders bei Regelvorschlägen dürfen fehlende Parameter nicht
durch erfundene Angaben ergänzt werden. Wenn beispielsweise eine
Zieladresse nicht eindeutig bestimmt werden kann, muss der
Vorschlag als unvollständig gekennzeichnet und zur weiteren
Prüfung zurückgehalten werden.

Für die dauerhafte Speicherung wird PostgreSQL als relationale
Datenbank vorgesehen. Darin können Sicherheitsereignisse,
Analysen, Regelvorschläge, Prüfresultate und Entscheidungen
miteinander verknüpft werden. Eine eindeutige Zuordnung zwischen
Ereignis und Vorschlag ermöglicht es, die Grundlage einer
Entscheidung später nachzuvollziehen.

\footnote{Vgl. PostgreSQL, \url{https://www.postgresql.org/docs/};
Ollama, \url{https://docs.ollama.com/capabilities/structured-outputs}.}

\subsection{Sicherheitskonzept und kontrollierte Freigabe}

Ein wesentliches Architekturprinzip ist die Trennung zwischen
der Erstellung eines KI-Vorschlags und seiner tatsächlichen
Ausführung. Das Sprachmodell soll ausschließlich Analysen und
Empfehlungen erzeugen. Es erhält keine direkten Berechtigungen
zur eigenständigen Änderung produktiver Sicherheitskonfigurationen.

Nach der Erstellung wird ein Regelvorschlag durch das Backend
auf strukturelle Vollständigkeit und technische Plausibilität
geprüft. Anschließend wird er einem kontrollierten Prüfprozess
zugeführt. Für die Erprobung von Änderungen ist eine isolierte
Testumgebung vorgesehen, in der mögliche Auswirkungen untersucht
werden können, ohne unmittelbar produktive Systeme zu verändern.

Erst nach erfolgreicher Prüfung und ausdrücklicher Freigabe durch
eine autorisierte Person darf eine separate Ausführungskomponente
die Änderung aktivieren. Die entsprechenden Berechtigungen und
Statusübergänge müssen serverseitig durchgesetzt werden. Eine
bloße Einschränkung der Schaltflächen im Frontend ist dafür
nicht ausreichend.

Die einzelnen Schritte sollen nachvollziehbar protokolliert
werden. Dazu gehören der ursprüngliche Vorschlag, die
Validierungsergebnisse, die Testresultate, die Freigabeentscheidung
und eine gegebenenfalls erfolgte Aktivierung. Zusätzlich ist
eine Möglichkeit zur Rücknahme problematischer Änderungen
vorzusehen.

Durch diese Aufteilung bleiben die Funktionen des Sprachmodells
von den sicherheitskritischen Operationen getrennt. Die KI
unterstützt den Analysten bei der Beurteilung eines Vorfalls,
während die Anwendung die Validierung, Autorisierung und
kontrollierte Ausführung übernimmt. Die konkrete technische
Umsetzung und Überprüfung dieser Abläufe erfolgt im Rahmen
der praktischen Realisierung.
